% chapters/preface.tex — はじめに

\chapter*{ヘッダー}
\addcontentsline{toc}{chapter} {ヘッダー}
\markboth{ヘッダー}{ヘッダー}

\section*{なぜこの本を書いたのか}

2022年末にChatGPTが公開された後、大型言語モデル（LLM）は人工知能分野の最も熱いテーマになった。数多くの開発者や研究者が、HuggingFace Transformersを介してすでに作成されたモデルをダウンロードし、ファインチューニングし、サービスに統合する方法を迅速に習得しました。しかし、「このモデルが内部的にどのように機能するのか」という質問に自信を持って答えることができる人はまだ少数です。

この本はその質問に答えるために使われました。トークナイザーを直接作成し、アテンションを手で実装し、トランスフォーマブロックを組み立て、学習ループを振り返りながらLLMのスケルトンを理解することがこの本の目標だ。抽象化されたライブラリを書かずにPyTorchの基本演算だけで最初から最後まで直接作る。その過程で経験するすべてのディテールが最終的にLLMを「理解」する道です。

\section*{この本の特徴}

この本は単に理論をリストしていません。\textbf{コードを最初に作成して, ユニットテストで検証した後, その過程を本に解く.}本に登場するすべてのPythonコードは実際に実行され、pytestベースのユニットテストに合格したコードのみが収録されました。読者はGitHubリポジトリからコード全体をダウンロードして直接実行、修正、学習することができます。

各章は、前の章で実装したモジュールをインポートして使用します。たとえば、4枚の埋め込みモジュールは3枚のトルクナイザーを使用し、5枚のアテンションは4枚の埋め込みを使用します。この「積み重ねられた構造」により、読者はLLMがどのように組み立てられるかを段階的に体感することができます。数学式は直感的な説明の後に登場し、初心者でも流れを逃さないように配慮した。

\section*{誰が読むべきか}

この本の主読者はPythonの基礎文法を知っていますが、ディープラーニングは初めての開発者です。 HuggingFace Transformersを書いたことがあるが、内部の動作を知らないエンジニア、「Attention Is All You Need」の論文を読んだが、コードで実装したことはない研究者も含まれる。微積分と基礎線形代数を仮定しますが、行列微分などの高度なトピックについては付録で補足説明します。

\begin{flushright}
\textit{2026年夏}\\[0.3em]
\textbf{dirmich}
\end{flushright}

\newpage

\section*{ありがとうの記事}

この本を書いている間、多くのオープンソースプロジェクトと論文の助けを受けました。特にAndrej KarpathyのnanoGPT、HuggingFace Tokenizers、そしてPyTorch開発チームの優れた文書に深く感謝します。彼らはLLMの民主化に専念した先駆者です。

また、この本のコードが基づいている多くの研究者に感謝します。 Ashish Vaswaniとトランスフォーマー原論文の共著者は、これらすべての出発点を作りました。 Alec RadfordとOpenAIチームは、GPTシリーズを通じてスケーリングの可能性を示した。 Tianyi ZhangとJason Weiは、インストラクションチューニングとソートの基礎を磨きました。

家族や友人にも感謝を伝える。執筆中に黙々と待ってくれた人々の忍耐がなければ、この本は完成しなかっただろう。

最後に、この本を読むすべての読者に。あなたの好奇心と情熱がLLMの未来を築きます。自分で作り、実験し、失敗し、もう一度やり直してください。それが学習の本質です。

\begin{flushright}
\textbf{dirmich}\\
\textit{Highmaru Press}
\end{flushright}
